\title{
Protokoll från Västra Götalandsregionen
}
\section*{Protokoll från beredningen för nära vård den 8 februari 2024}
Tid: 09:30 - 11:00
Plats: Lokal Bommen, Stationshuset i Göteborg
\section*{Mötesuppehåll}
Uppehåll kl. 10:25-10:35
\section*{Närvarande}
\section*{Beslutande}
Janette Olsson, ordförande (S)
Pär Lundqvist, 1:a vice ordförande (L)
My Alnebratt (S)
Tomas Johansson (M)
Håkan Lösnitz (SD)
Carina Örgård (V)
Per-Olof Blomqvist (KD)
Cecilia Andersson (C)
\section*{Justerare}
Pär Lundqvist (L)
\section*{Justering}
Protokollet justeras elektroniskt.
\section*{Sekreterare}
Katarina Trampush
\section*{Underskrifter}
\section*{Ordförande:}
Janette Olsson (S)
\begin{tabular}{lllll}
\hline \begin{tabular}{l} 
Potadress: \\
Regionens hus \\
46280 Vänersborg
\end{tabular} & \begin{tabular}{l} 
Beskvaldress: \\
Residensjatan 16H \\
46233 Vänersborg
\end{tabular} & \begin{tabular}{l}
Telefon: \\
1016-441 00 00
\end{tabular} & Webplats: & E-post: \\
\hline
\end{tabular}

Justerare:
Pär Lundqvist (L)

\title{
Politiska sekreterare
}
AnnaSofia Alexandersson (MP), distans
Johan Sösaeter Johansson (SD)
Mats Berglund (V)
Moez Dharsani (S)
Nanna Siewerts Tulinius (L)
Emma Askengren (C), § 1-3, distans
\section*{Övriga närvarande}
Karin Looström Muth, hälso- och sjukvårdsutvecklingsdirektör
Kaarina Sundelin, produktionsdirektör hälso- och sjukvård
Ann-Sofi Isaksson, avdelningschef, primärvård och regionövergripande
verksamheter
Annika Mårtensson, ledningsstöd till hälso- och sjukvårdsutvecklingsdirektör
Lillemor Harnell, ledningsstöd, koncernstab strategisk hälso- och
sjukvårdsutveckling
Katarina Trampush, nämndsamordnare, avdelning nämndsamordning
Matteus Berglund, nämndsamordnare, avdelning nämndsamordning
Nina Bonnedahl, enhetschef ledningsstöd, § 2-3
Anna Aronsson, regionutvecklare, primärvård och regionövergripande
verksamheter, § 4, distans

Protokoll från beredningen för nära vård, 2024-02-08
\title{
Innehållsförteckning
}
\(\S 1 \quad\) Information om remiss strategisk inriktning för vårdvalssystemet i Västra 5
Götalandsregionen 2025-2028
\(\S 2 \quad\) Lägesrapport i arbetet med avtal och uppdrag
\(\S 3 \quad\) Lägesrapport i arbetet med ekonomiska principer och ersättningsmodeller
\(\S 4 \quad\) Svar på SSN uppdrag: att bedöma om det finns behov av ytterligare beslut för att den operativa hälso- och sjukvårdsnämnden ska kunna agera på kort sikt i arbetet mot nationella riktvärdet för fast läkarkontakt i primärvärden

Protokoll från beredningen för nära vård, 2024-02-08
\section*{\(\S 1\)}
\section*{Information om remiss strategisk inriktning för vårdvalssystemet i Västra Götalandsregionen 2025-2028}
Diarienummer SSN 2023-00656
\section*{Beslut}
1. Beredningen för nära vård noterar informationen.
2. Beredningen för nära vård ställer sig bakom att skicka ut remiss om strategisk inriktning för vårdvalssystemet i Västra Götalandsregionen 202 2028.
\section*{Sammanfattning av ärendet}
Ett förslag till inriktning för vårdvalssystemet i Västra Götalandsregionen har tagits fram av beredningen för nära vård, utifrån beredningens uppdrag att ta fram underlag för inriktning inför ändring av vårdvalssystemet inför 2025 som operativa hälso- och sjukvårdsnämnden sedan ska kunna omsätta till förfrågningunderlag.
Den strategiska inriktningen utgår från beslutad strategi och genomförandeplan för omställningen samt insamlade krav och önskemål från beredningen för nära vård. Som en del i arbetet med inriktningen har även dialog genomförts med och synpunkter inhämtats från facklig referensgrupp och patienter, genom Levande bibliotek.
Insamlade krav och önskemål har bearbetats och mynnat ut i ett förslag till strategisk inriktning med fem prioriterade områden:
- Proaktivt genom hälsofrämjande och förebyggande insatser
- Stärkt samverkan och samordning
- Ökad tillgänglighet och personcentrering
- Högsta kvalitet och bästa möjliga kunskap
- Prioritering och differentiering
Det finns även ett antal grundläggande förutsättningar som en del av inriktningen; kultur och ledarskap, kompetensförsjörjning, digitalisering som verktyg samt ekonomiskt ansvarstagande.
Under sammanträdet får beredningen information om förslag till strategisk inriktning, som kommer gå ut på remiss till berörda nämnder och aktörer under

Protokoll från beredningen för nära vård, 2024-02-08
februari 2024. Remisstiden är tre månader och som remissmottagare står operativa hälso- och sjukvårdsnämnden, samtliga delregionala nämnder, samtliga utförarstyrelser inom hälso- och sjukvård samt privata vårdgivare inom vårdval vårdcentral och vårdval rehab genom PrimÖR. Synpunkter från det politiska samrådsorganet med länets kommuner (SRO) kommer att hanteras genom workshop den 6 mars 2024.
Ärendet är ett informationsärende och medför inga finansierings- eller resurskonsekvenser.
Beredningen för nära vård enas under sammanträdet om att ställa sig bakom att skicka ut remiss om strategisk inriktning för vårdvalssystemet i Västra
Götalandsregionen 2025-2028.
\section*{Beslutsunderlag}
- Tjänsteutlätande daterat 2023-12-19